% =====================================================================
%  Lagrange Duality: Dual Function Construction, Strong Duality,
%  and Sensitivity Analysis
%  Amila Snagic --- Numerical Optimization
%
%  Compiles with pdflatex (run twice for the table of contents and
%  cross-references).  Figures are expected in ../figures or ./figures.
% =====================================================================
\documentclass[11pt,a4paper]{article}

\usepackage[T1]{fontenc}
\usepackage[utf8]{inputenc}
\usepackage[a4paper,top=2.3cm,bottom=2.3cm,left=2.4cm,right=2.4cm]{geometry}

\usepackage{amsmath}
\usepackage{amssymb}
\usepackage{amsthm}
\usepackage{graphicx}
\usepackage{booktabs}
\usepackage{xcolor}
\usepackage{listings}
\usepackage{helvet}          % Helvetica-style sans, used on the title page
\usepackage[hidelinks]{hyperref}

\graphicspath{{../figures/}{figures/}{./}}

% ---------------------------------------------------------------- theorem styles
\theoremstyle{definition}
\newtheorem{proposition}{Proposition}
\newtheorem{theorem}{Theorem}

% ---------------------------------------------------------------- code listings
\definecolor{codebg}{RGB}{245,245,245}
\definecolor{codekw}{RGB}{0,0,160}
\definecolor{codecm}{RGB}{100,100,100}
\lstset{
  language=C++,
  basicstyle=\ttfamily\footnotesize,
  keywordstyle=\color{codekw},
  commentstyle=\color{codecm}\itshape,
  backgroundcolor=\color{codebg},
  frame=none,
  breaklines=true,
  showstringspaces=false,
  xleftmargin=8pt,
  framexleftmargin=8pt,
  aboveskip=8pt,
  belowskip=8pt,
}

% ---------------------------------------------------------------- shorthands
\newcommand{\R}{\mathbb{R}}
\newcommand{\pstar}{p^{\star}}
\newcommand{\dstar}{d^{\star}}
\newcommand{\xstar}{x^{\star}}
\newcommand{\lamstar}{\lambda^{\star}}
\newcommand{\conv}{\operatorname{\mathbf{conv}}}
\DeclareMathOperator*{\argmin}{arg\,min}

\setlength{\parskip}{0.45em}
\setlength{\parindent}{0pt}

\begin{document}

% ================================================================= title page
\begin{titlepage}
\centering
\vspace*{4.2cm}

{\sffamily\bfseries\fontsize{24}{30}\selectfont
Lagrange Duality: Dual Function\\[2pt]
Construction, Strong Duality, and\\[2pt]
Sensitivity Analysis\par}

\vspace{2.0cm}
{\sffamily\fontsize{14}{17}\selectfont Numerical Optimization\par}

\vspace{1.0cm}
{\sffamily\fontsize{12.5}{15}\selectfont
Interactive implementation on the natID framework\par}

\vspace{1.7cm}
{\sffamily\fontsize{14}{17}\selectfont Amila Snagic\par}

\vspace{1.7cm}
{\sffamily\fontsize{12}{14}\selectfont
University of Sarajevo --- Faculty of Electrical Engineering\par}

\vspace{0.6cm}
{\sffamily\itshape\fontsize{11.5}{14}\selectfont Academic year 2025/26\par}

\vfill
\end{titlepage}

\setcounter{page}{1}

% ================================================================= abstract
\begin{center}
\fcolorbox{gray!40}{gray!5}{%
\begin{minipage}{0.95\textwidth}
\small
\textbf{Abstract.} The Lagrange dual function is constructed by pointwise minimisation of the
Lagrangian, and weak duality, strong duality, Slater's condition and the shadow-price reading of
the multipliers are derived. Three instances are implemented: a strictly convex quadratic program
whose dual is available in closed form, a resource-allocation problem reproducing the
water-filling rule, and a non-convex problem where strong duality fails. In the convex cases the
primal is solved twice by unrelated methods, so the reported duality gap is a cross-validation
rather than one number printed twice. For the non-convex instance the perturbation function and
its convex envelope exhibit the gap directly. All results agree with theory to between
$10^{-11}$ and $10^{-16}$.
\end{minipage}}
\end{center}

\vspace{0.6cm}
\tableofcontents
\newpage

% ================================================================= 1
\section{Introduction}

Duality attaches to every constrained minimisation problem a second, always-convex problem built
from its constraints. The dual's value bounds the primal's from below, under mild conditions
equals it, and its solution prices the constraints. This report derives that theory, applies it
to three instances, and describes the interactive application built to compute and draw it.
Every figure was produced by that application.

% ================================================================= 2
\section{Mathematical foundations}\label{sec:theory}

\subsection{The Lagrangian and the dual function}

For the problem in standard form
\begin{equation}\label{eq:primal}
\min_{x\in\R^{n}} f_{0}(x)\quad\text{subject to}\quad f_{i}(x)\le 0,\quad i=1,\dots,m,
\end{equation}
with optimal value $\pstar$, the \textbf{Lagrangian} attaches a price $\lambda_{i}$ to each
constraint,
\begin{equation}
L(x,\lambda)=f_{0}(x)+\sum_{i=1}^{m}\lambda_{i}f_{i}(x),
\end{equation}
and the \textbf{dual function} is its unconstrained infimum over $x$,
\begin{equation}\label{eq:dual}
g(\lambda)=\inf_{x\in\R^{n}}L(x,\lambda),
\end{equation}
which may equal $-\infty$. The \textbf{dual problem} is $\max_{\lambda\ge 0}g(\lambda)$, with
optimal value $\dstar$. No convexity is assumed anywhere above.

\begin{proposition}\label{prop:concave}
For every problem of the form \eqref{eq:primal} --- convex or not --- $g$ is concave.
\end{proposition}

\begin{proof}
For fixed $x$, the map $\lambda\mapsto L(x,\lambda)$ is \emph{affine}, since $f_{0}(x)$ and
$f_{i}(x)$ are then constants. By \eqref{eq:dual}, $g$ is a pointwise infimum of affine
functions, and such an infimum is concave.
\end{proof}

This is the structural reason duality is useful: whatever pathology the primal has,
\textbf{the dual is a concave maximisation over a convex set}. Section~\ref{sec:gap} exhibits a
non-convex primal with a perfectly well-behaved concave dual.

\subsection{Weak duality}

\begin{theorem}[Weak duality]\label{thm:weak}
For every $\lambda\ge 0$, $g(\lambda)\le\pstar$; hence $\dstar\le\pstar$.
\end{theorem}

\begin{proof}
Let $\tilde{x}$ be feasible and $\lambda\ge 0$. Each $\lambda_{i}f_{i}(\tilde{x})$ is
non-positive, so
\[
g(\lambda)=\inf_{x}L(x,\lambda)\le L(\tilde{x},\lambda)
=f_{0}(\tilde{x})+\sum_{i}\lambda_{i}f_{i}(\tilde{x})\le f_{0}(\tilde{x}).
\]
Taking the infimum over feasible $\tilde{x}$ and the supremum over $\lambda\ge 0$ gives the
claim.
\end{proof}

Only the sign conditions $\lambda\ge 0$ and $f_{i}(\tilde{x})\le 0$ were used ---
\textbf{no convexity}. Weak duality is therefore universal, which is what makes dual bounds
valuable in non-convex optimisation. The quantity $\pstar-\dstar\ge 0$ is the
\textbf{duality gap}.

\subsection{Strong duality, Slater's condition, and the KKT system}

Equality is not automatic. \textbf{Slater's condition} is the standard sufficient criterion: if
the problem is convex and there exists a \textbf{strictly feasible} point
($f_{i}(\tilde{x})<0$ for all $i$), then $\dstar=\pstar$ and the dual optimum is attained. It is
sufficient, not necessary.

When strong duality holds and both optima are attained at $\xstar,\lamstar$, the chain
\begin{equation}
f_{0}(\xstar)=g(\lamstar)\le f_{0}(\xstar)+\sum_{i}\lambda_{i}^{\star}f_{i}(\xstar)\le f_{0}(\xstar)
\end{equation}
begins and ends at the same value, so both inequalities are equalities. The second forces every
term of $\sum_{i}\lambda_{i}^{\star}f_{i}(\xstar)$ --- each $\le 0$ --- to vanish:
\begin{equation}\label{eq:cs}
\lambda_{i}^{\star}f_{i}(\xstar)=0,\qquad i=1,\dots,m.
\end{equation}
This is \textbf{complementary slackness}: a constraint is either active and may carry any
non-negative price, or slack and priced at zero. The first equality says $\xstar$ minimises
$L(\cdot,\lamstar)$ over all of $\R^{n}$, so in the differentiable case
\begin{equation}\label{eq:stat}
\nabla f_{0}(\xstar)+\sum_{i}\lambda_{i}^{\star}\nabla f_{i}(\xstar)=0
\end{equation}
(\textbf{stationarity}). With primal and dual feasibility these are the \textbf{KKT conditions},
necessary and sufficient for a convex problem satisfying Slater. All four residuals are measured
in Section~\ref{sec:verify}.

\subsection{Sensitivity: shadow prices and the convex envelope}\label{sec:pert}

Perturbing the constraint levels defines the \textbf{perturbation function}
\begin{equation}\label{eq:pu}
p(u)=\inf\{f_{0}(x):f_{i}(x)\le u_{i}\},\qquad p(0)=\pstar.
\end{equation}

\begin{theorem}[Global sensitivity]\label{thm:sens}
If strong duality holds with dual optimum $\lamstar$, then
$p(u)\ge\pstar-(\lamstar)^{T}u$ for every $u$.
\end{theorem}

\begin{proof}
If $f_{i}(x)\le u_{i}$ then, using $\lamstar\ge 0$,
$\pstar=g(\lamstar)\le f_{0}(x)+\sum_{i}\lambda_{i}^{\star}f_{i}(x)\le f_{0}(x)+(\lamstar)^{T}u$.
Rearranging and taking the infimum gives the claim.
\end{proof}

So $u\mapsto\pstar-(\lamstar)^{T}u$ is a supporting hyperplane touching $p$ at the origin, and
$-\lamstar$ is a subgradient of $p$ there. Where $p$ is differentiable,
\begin{equation}\label{eq:shadow}
\frac{\partial\pstar}{\partial u_{i}}=-\lambda_{i}^{\star}:
\end{equation}
$\lambda_{i}^{\star}$ is the \textbf{shadow price} of constraint $i$, the rate at which the
optimal cost falls per unit of relaxation. An inactive constraint has $\lambda_{i}^{\star}=0$ by
\eqref{eq:cs} and so prices at zero, as intuition demands.

The same perturbation view explains gaps. One shows
$g(\lambda)=\inf_{u}\bigl(p(u)+\lambda^{T}u\bigr)$, whence $\dstar$ is the value at $u=0$ of the
largest convex function lying below $p$ --- its \textbf{convex envelope}:
\begin{equation}\label{eq:gapgeom}
\pstar=p(0),\qquad \dstar=(\conv p)(0),\qquad \text{gap}=p(0)-(\conv p)(0).
\end{equation}
If $p$ is convex near the origin the two coincide. If $p$ has a non-convex dent there, the
envelope cuts a chord across it and \textbf{the gap is the depth of the dent} --- computed
directly in Section~\ref{sec:gap}.

% ================================================================= 3
\section{Instance A: convex quadratic program}\label{sec:qp}

\subsection{Model and closed-form dual}

\begin{equation}\label{eq:qp}
\min_{x\in\R^{n}}\ \tfrac{1}{2}x^{T}Px+q^{T}x\quad\text{subject to}\quad Ax\le b,
\qquad P\succ 0,
\end{equation}
with default data
\begin{equation}\label{eq:qpdata}
P=\begin{pmatrix}2&1\\1&2\end{pmatrix},\quad
q=\begin{pmatrix}-3\\-4\end{pmatrix},\quad
A=\begin{pmatrix}1&1\\-1&2\end{pmatrix},\quad
b=\begin{pmatrix}2\\4\end{pmatrix},
\end{equation}
chosen so one constraint binds at the optimum and the other does not, giving complementary
slackness something visible to say.

$L$ is strictly convex in $x$ for every $\lambda$ (its Hessian is $P\succ 0$), so the infimum in
\eqref{eq:dual} is attained at the stationary point:
\begin{align}
Px+q+A^{T}\lambda=0 &\quad\Longrightarrow\quad x(\lambda)=-P^{-1}(q+A^{T}\lambda),
\label{eq:xlam}\\[4pt]
g(\lambda)&=-\tfrac{1}{2}(q+A^{T}\lambda)^{T}P^{-1}(q+A^{T}\lambda)-b^{T}\lambda.
\label{eq:gqp}
\end{align}
Discarding a constant, maximising $g$ over $\lambda\ge 0$ is the \textbf{non-negatively
constrained QP}
\begin{equation}\label{eq:dualqp}
\min_{\lambda\ge 0}\ \tfrac{1}{2}\lambda^{T}H\lambda+d^{T}\lambda,\qquad
H=AP^{-1}A^{T}\succeq 0,\qquad d=b+AP^{-1}q.
\end{equation}
The dual of a QP is a QP, but over the far simpler set $\lambda\ge 0$ --- the practical payoff of
duality here.

\subsection{Solving the dual: exact cyclic coordinate ascent}

Fixing all coordinates but the $i$-th, \eqref{eq:dualqp} restricts to the scalar parabola
$\tfrac{1}{2}H_{ii}\lambda_{i}^{2}+\lambda_{i}\bigl(d_{i}+\sum_{k\ne i}H_{ik}\lambda_{k}\bigr)$.
Its minimiser over the half-line $\lambda_{i}\ge 0$ is the vertex, projected:
\begin{equation}\label{eq:cd}
\lambda_{i}\ \leftarrow\ \max\!\left(0,\
-\frac{d_{i}+\sum_{k\ne i}H_{ik}\lambda_{k}}{H_{ii}}\right).
\end{equation}
Each update is exact and in closed form --- no step size, no line search. Sweeps run until the
largest coordinate change falls below $10^{-12}$. Because the feasible set is a product of
half-lines, coordinate-wise minimisation cannot stall at a non-stationary point, the difficulty
that makes coordinate descent unreliable for non-separable constraints. Coordinates with
$H_{ii}\approx 0$ are skipped, handling degenerate constraint rows. The primal is then recovered
from \eqref{eq:xlam}, which \emph{is} the stationarity condition, so \eqref{eq:stat} holds to
machine precision automatically.

\subsection{Solving the primal independently: logarithmic barrier method}

Recovering $\xstar$ from $\lamstar$ would make $\pstar$ and $\dstar$ two readings of one
computation, so the gap would prove nothing. The primal is therefore solved a second time by an
unrelated route. For $t>0$ define
\begin{equation}\label{eq:barrier}
\psi_{t}(x)=t\left(\tfrac{1}{2}x^{T}Px+q^{T}x\right)-\sum_{i=1}^{m}\log(b_{i}-a_{i}^{T}x),
\end{equation}
whose barrier term is finite exactly on $Ax<b$ and diverges at the boundary. Its gradient and
Hessian are $\sum_{i}a_{i}/(b_{i}-a_{i}^{T}x)$ and
$\sum_{i}a_{i}a_{i}^{T}/(b_{i}-a_{i}^{T}x)^{2}\succeq 0$, so $\psi_{t}$ is strictly convex with a
unique minimiser $x^{\star}(t)$; the curve $t\mapsto x^{\star}(t)$ is the \textbf{central path},
and $f_{0}(x^{\star}(t))-\pstar\le m/t$.

The implementation uses the classical two loops. The outer loop starts at $t=1$ and multiplies by
$\mu=12$ until $m/t<10^{-10}$. The inner loop applies Newton's method, solving
$\nabla^{2}\psi_{t}\,\Delta x=-\nabla\psi_{t}$ by LU with partial pivoting and stopping on the
\textbf{Newton decrement} $\lambda_{N}^{2}=-\nabla\psi_{t}^{T}\Delta x$ when
$\lambda_{N}^{2}/2\le 10^{-12}$; the decrement is used because it is affine-invariant and
estimates $\psi_{t}(x)-\inf\psi_{t}$ directly. Backtracking ($\alpha=0.25$, $\beta=0.5$) accepts
a step only when it is \emph{strictly feasible} and satisfies the Armijo condition, which is what
keeps iterates in the domain of the logarithms.

A strictly feasible start is produced by minimising the smooth squared-hinge measure
$\Psi(x)=\sum_{i}\max(0,\,a_{i}^{T}x-b_{i}+\delta)^{2}$ with $\delta=10^{-3}$, by gradient
descent with an adaptive step. Squaring removes the kink of $\max(0,\cdot)$, and $\Psi(x)=0$
certifies strict feasibility with margin $\delta$ --- so the same routine both initialises the
barrier method and answers Slater's condition constructively.

\subsection{Results, complementary slackness, and sensitivity}

The two solvers return
\[
\xstar=(0.5,\ 1.5)^{T},\qquad \lamstar=(0.5,\ 0)^{T},\qquad
\pstar=\dstar=-4.25,\qquad \text{gap}=0.
\]
The barrier method, which never touches the dual, returns $-4.249999999984$ --- agreement to
$1.6\times 10^{-11}$, which is what makes the zero gap meaningful. Complementary slackness is
exact: constraint~1 binds ($\lambda_{1}^{\star}=0.5$, slack $0$), constraint~2 is slack
($\lambda_{2}^{\star}=0$, slack $1.5$), and both products vanish, exactly as \eqref{eq:cs}
requires.

Figure~\ref{fig:primal} shows the geometry: objective contours, the feasible polyhedron, the
unconstrained minimiser $-P^{-1}q=(2/3,\,5/3)$ lying outside it, and $\xstar$ where the innermost
reachable contour touches the active constraint --- the geometric content of \eqref{eq:stat}.
Figure~\ref{fig:dual} plots $g$ over the non-negative quadrant: a single concave hill whose
summit lies \emph{on the boundary} $\lambda_{2}=0$, the dual-side image of constraint~2 being
inactive.

\begin{figure}[htbp]
\centering
\includegraphics[width=0.62\textwidth]{1_primal_qp.pdf}
\caption{Instance A, primal geometry. Filled bands are level sets of $f_{0}$, the hatched wedge
is $\{x:Ax\le b\}$, the active constraint is solid red and the inactive one dashed.}
\label{fig:primal}
\end{figure}

\begin{figure}[htbp]
\centering
\includegraphics[width=\textwidth]{2_dual_function.pdf}
\caption{Instance A, the dual function $g(\lambda_{1},\lambda_{2})$ on $\lambda\ge 0$. Warm
colours are high. Every point is a lower bound on $\pstar$ (Theorem~\ref{thm:weak}); the summit
attains it.}
\label{fig:dual}
\end{figure}

Figure~\ref{fig:sens} perturbs the first constraint to $a_{1}^{T}x\le b_{1}+u$, re-solving at
each $u$, and compares with the tangent of Theorem~\ref{thm:sens}. The numerical derivatives
match the multipliers: $\partial\pstar/\partial u_{1}=-0.5=-\lambda_{1}^{\star}$ to
$5.6\times 10^{-12}$, and $\partial\pstar/\partial u_{2}=0=-\lambda_{2}^{\star}$. The curve is
equally instructive where it \emph{stops} matching: for $u>0$ it flattens, because relaxing an
already-relaxed constraint eventually makes it inactive. That is the active set changing, and it
is exactly where $p$ loses differentiability and the tangent only bounds, as
Theorem~\ref{thm:sens} says it must.

\begin{figure}[htbp]
\centering
\includegraphics[width=\textwidth]{3_sensitivity.pdf}
\caption{Instance A, sensitivity. Solid: $p^{\star}(u)$ recomputed exactly. Dashed: the
prediction $\pstar-\lambda_{1}^{\star}u$. Tangency at the origin is \eqref{eq:shadow}.}
\label{fig:sens}
\end{figure}

% ================================================================= 4
\section{Instance B: resource allocation and water filling}\label{sec:alloc}

\subsection{Model, separable dual, and the water-filling rule}

\begin{equation}\label{eq:alloc}
\max_{x\ge 0}\ \sum_{i=1}^{n}w_{i}\log\!\left(1+\frac{x_{i}}{\alpha_{i}}\right)
\quad\text{subject to}\quad \sum_{i=1}^{n}x_{i}\le B,
\end{equation}
with $w_{i},\alpha_{i}>0$. As a minimisation,
$f_{0}(x)=-\sum_{i}w_{i}\log(1+x_{i}/\alpha_{i})$, which is convex, subject to
$\sum_{i}x_{i}-B\le 0$.

\textbf{Only the budget is dualised}; the constraints $x\ge 0$ are kept in the domain of the
inner minimisation. This leaves one dual variable, so the dual function is a curve that can
simply be drawn --- and it is what makes the closed form below come out. The Lagrangian then
separates over $i$, and each scalar piece is convex with derivative
$-w_{i}/(\alpha_{i}+x_{i})+\lambda$, giving
\begin{equation}\label{eq:waterfill}
x_{i}(\lambda)=\max\!\left(0,\ \frac{w_{i}}{\lambda}-\alpha_{i}\right).
\end{equation}
This is the classical \textbf{water-filling} rule: $w_{i}/\lambda$ is a \emph{water level},
$\alpha_{i}$ the \emph{floor} of basin $i$, and $x_{i}$ the depth of water above that floor.
Basins whose floor already lies above their level receive nothing. For $\lambda\le 0$ the
Lagrangian is unbounded below, so the dual domain is $\lambda>0$.

Each $x_{i}(\lambda)$ is non-increasing in $\lambda$, so the total
$T(\lambda)=\sum_{i}x_{i}(\lambda)$ is continuous and strictly decreasing where positive, with
$T\to\infty$ as $\lambda\to 0^{+}$ and $T\to 0$ as $\lambda\to\infty$. Hence $T(\lambda)=B$ has a
\emph{unique} root, located by bisection: the bracket is found by doubling, then 300 halvings.
Bisection is the right tool despite its linear rate --- $T$ has kinks wherever a basin enters or
leaves the water, which would need care with Newton's method, whereas bisection needs only
monotonicity and continuity.

At the root the budget binds, and the KKT conditions read
\begin{equation}\label{eq:market}
\frac{w_{i}}{\alpha_{i}+x_{i}^{\star}}=\lamstar \ \text{ if } x_{i}^{\star}>0,
\qquad\qquad
\frac{w_{i}}{\alpha_{i}}\le\lamstar \ \text{ if } x_{i}^{\star}=0,
\end{equation}
the market-clearing reading: at the optimum all funded activities are \textbf{equally} worth
funding, and anything unfunded is not worth the price of the budget even at zero allocation.

\subsection{Results}

With $w=(3,2,1,0.5)$, $\alpha=(1,0.5,2,0.2)$, $B=6$:

\begin{center}
\begin{tabular}{ccccl}
\toprule
$i$ & $w_{i}$ & $\alpha_{i}$ & $x_{i}^{\star}$ &
\ \ marginal utility $w_{i}/(\alpha_{i}+x_{i}^{\star})$\\
\midrule
1 & 3.0 & 1.0 & 3.2 & \ \ 0.714285714\\
2 & 2.0 & 0.5 & 2.3 & \ \ 0.714285714\\
3 & 1.0 & 2.0 & 0.0 & \ \ 0.500000000 --- shut out\\
4 & 0.5 & 0.2 & 0.5 & \ \ 0.714285714\\
\bottomrule
\end{tabular}
\end{center}

with $\lamstar=0.714285714$, utility $8.3771683$ and $\pstar=\dstar$ to twelve digits. The table
is \eqref{eq:market} made concrete: the three funded activities share the common marginal utility
$\lamstar$, while activity 3 --- whose floor is high relative to its weight --- sits at $0.5$,
below the price, and is correctly shut out. The shadow price is confirmed numerically,
$\partial\pstar/\partial B=-\lamstar$ to $2.4\times 10^{-11}$: one extra unit of budget buys
$\lamstar$ extra utility.

\begin{figure}[htbp]
\centering
\includegraphics[width=0.60\textwidth]{4b_water_filling.pdf}
\caption{Instance B. Grey bars are floors $\alpha_{i}$, blue the allocation $x_{i}$, red segments
the levels $w_{i}/\lamstar$. Basin 3's floor is above its level, so it stays dry.}
\label{fig:water}
\end{figure}

% ================================================================= 5
\section{Instance C: a non-convex problem with a duality gap}\label{sec:gap}

\subsection{Model and method}

\begin{equation}\label{eq:nc}
\min_{x\in\R} f_{0}(x)=x^{4}+c_{2}x^{2}+c_{1}x\quad\text{subject to}\quad x\ge x_{\min},
\end{equation}
with $c_{2}=-6$, $c_{1}=1$, $x_{\min}=0$. For $c_{2}<0$ this is a \textbf{double well}; with
these defaults the global minimum lies at $x\approx-1.7589$ with value $\approx-10.7524$, and
\textbf{the constraint excludes it}. Writing the constraint as $f_{1}(x)=x_{\min}-x\le 0$,
\begin{equation}\label{eq:gnc}
g(\lambda)=\lambda x_{\min}+\inf_{x}\bigl[x^{4}+c_{2}x^{2}+(c_{1}-\lambda)x\bigr],
\end{equation}
so \textbf{dualising merely shifts the linear coefficient}, $c_{1}\mapsto c_{1}-\lambda$, and one
routine evaluates both $f_{0}$ and $g$. Since $f_{0}$ is coercive the infimum is finite for every
$\lambda$, so $g$ is real-valued --- and concave by Proposition~\ref{prop:concave},
\textbf{even though $f_{0}$ is not convex}.

Both $f_{0}$ and $g$ need \emph{global} minimisation, so the implementation uses a 4001-point
grid scan followed by 120 golden-section iterations within the two cells adjacent to the best
grid point. For a coercive quartic with at most three critical points a grid this fine cannot
miss the global basin, and the refined bracket is unimodal by construction. The outer
maximisation over $\lambda$ uses the same scheme, legitimately, because $g$ is concave. The
convex envelope of the sampled $p(u)$ is computed as a \textbf{lower convex hull} by a
monotone-chain sweep in $O(N)$, popping a point whenever it lies on or above the chord joining
its neighbours.

\subsection{Results: the gap, three ways}

\[
\xstar=1.6887758,\quad \pstar=-7.2893095,\quad \lamstar=1.0,\quad \dstar=-9.0,\quad
\text{gap}=1.7106905.
\]
$\dstar=-9$ is exact and verifiable by hand: at $\lambda=1$ the shifted linear coefficient
vanishes, so \eqref{eq:gnc} reduces to $\inf_{x}(x^{4}-6x^{2})=-9$ at $x=\pm\sqrt{3}$. That the
maximum falls there is the subgradient condition --- $g'(\lambda)=x_{\min}-x(\lambda)$, and at
$\lambda=1$ the inner problem has \emph{two} minimisers, so the subdifferential is
$[-\sqrt{3},\sqrt{3}]\ni 0$. The kink at the summit of Figure~\ref{fig:ncdual} is exactly this
non-uniqueness.

\begin{center}
\begin{tabular}{lr}
\toprule
\textbf{Route} & \textbf{Value at $u=0$}\\
\midrule
Primal optimum $\pstar=p(0)$ & $-7.2893094547$\\
Dual optimum $\dstar=\max_{\lambda\ge 0}g(\lambda)$ & $-9.0000000000$\\
Convex envelope $(\conv p)(0)$, from the hull & $-8.9992382400$\\
\bottomrule
\end{tabular}
\end{center}

The envelope value matches $\dstar$ to $7.6\times 10^{-4}$, the grid resolution --- a direct
numerical confirmation of \eqref{eq:gapgeom}.

\begin{figure}[htbp]
\centering
\includegraphics[width=0.47\textwidth]{5a_nonconvex_objective.pdf}
\caption{The objective: the global minimum at $x\approx-1.76$ is infeasible, and the constrained
optimum sits in the right-hand well.}
\label{fig:ncobj}
\end{figure}

\begin{figure}[htbp]
\centering
\includegraphics[width=0.74\textwidth]{5b_perturbation_and_envelope.pdf}
\caption{$p(u)$ and its convex envelope. $p$ is flat at $-7.289$, then drops to $-10.75$ near
$u\approx 1.7$ when the relaxation finally admits the deeper well. That shape is not convex; the
envelope cuts the chord across it, and the red segment at $u=0$ is the gap.}
\label{fig:pert}
\end{figure}

\begin{figure}[htbp]
\centering
\includegraphics[width=0.47\textwidth]{5c_nonconvex_dual.pdf}
\caption{The dual function: concave despite the non-convex primal, and staying strictly below
$\pstar$.}
\label{fig:ncdual}
\end{figure}

The instance separates the two theorems cleanly. \textbf{Weak duality survives} ---
$\dstar\le\pstar$ still holds and $g$ still certifies a lower bound, its proof having used only
sign conditions. \textbf{Strong duality fails} --- strictly feasible points exist in abundance
here (every $x>0$), so it is precisely the convexity hypothesis of Slater's condition that has
been removed, and the conclusion fails with it. The practical lesson is the standard one: in
non-convex optimisation the dual is still worth solving as a source of certified bounds, and
\eqref{eq:gapgeom} says in advance how weak those bounds can be.

% ================================================================= 6
\section{Implementation}\label{sec:impl}

\subsection{Architecture}

The application is about 4100 lines of C++20 on the natID framework, split so that the
mathematics is independent of the interface.

\begin{center}
\begin{tabular}{lrl}
\toprule
\textbf{File} & \textbf{Lines} & \textbf{Responsibility}\\
\midrule
\texttt{core/LinAlg.h}             & 219  & dense LU, Cholesky, vector primitives\\
\texttt{core/QPProblem.h}          & 402  & dual coordinate ascent, barrier primal, phase I\\
\texttt{core/AllocationProblem.h}  & 147  & water filling, bisection\\
\texttt{core/NonconvexProblem.h}   & 217  & grid $+$ golden section, convex envelope\\
\texttt{core/AppModel.h}           & 141  & the three instances, half-plane clipping\\
\texttt{views/PlotCanvas.h}        & 821  & plotting engine\\
\texttt{views/View*.h}             & 1550 & six tabs\\
\bottomrule
\end{tabular}
\end{center}

The \texttt{core/} directory has \textbf{no GUI dependency whatsoever}. This was deliberate and
had a practical payoff: the numerics were developed and verified by a standalone console program
compiled against \texttt{core/} alone, before any drawing code existed, and every number in
Section~\ref{sec:verify} comes from that harness. A single \texttt{AppModel} owns the three
instances; every parameter panel calls one shared callback that re-solves and repaints all six
tabs, so the application cannot display mutually inconsistent state.

\subsection{Numerical core}

\textbf{Cholesky} $P=LL^{T}$ is used wherever $P^{-1}$ appears. Beyond the usual factor-of-two
saving it is \textbf{backward stable without pivoting} for symmetric positive definite matrices,
and its failure is \textbf{diagnostic}: a non-positive radicand means $P$ is not positive
definite, so the same routine doubles as the positive-definiteness test reported in the
interface. \textbf{LU with partial pivoting} solves the Newton systems, where near the boundary
the barrier Hessian has entries of order $10^{10}$ and pivoting guards against the loss of
positive definiteness that rounding can induce. Forming $H=AP^{-1}A^{T}$ never inverts $P$: each
row $a_{i}$ is passed through the triangular solves, giving $m$ solves instead of an inversion,
both cheaper and better conditioned.

\begin{lstlisting}
static QPSolution solve(const QPData& qp)
{
    sol.spd = la::isSPD(qp.P);              // Cholesky, doubles as the PD test
    if (!sol.spd) return sol;               // dual is -infinity everywhere
    la::cholesky(qp.P, L);

    sol.lambda = solveDual(qp, L, sol.dualIters);   // coordinate ascent, dual QP
    xOfLambda(qp, L, sol.lambda, sol.x);            // x* = -P^-1 (q + A'lam)
    sol.pStar = f0(qp, sol.x);
    sol.dStar = dualValue(qp, L, sol.lambda);
    sol.gap   = sol.pStar - sol.dStar;

    sol.slater    = findStrictlyFeasible(qp, x0);   // phase I, also answers Slater
    sol.barrierOk = solvePrimalBarrier(qp, sol.xBarrier, sol.pStarBarrier);
    return sol;
}
\end{lstlisting}

For the default data the coordinate ascent converges in a \textbf{single sweep}: the first update
sets $\lambda_{1}$ exactly, the second finds $\lambda_{2}$'s unconstrained optimum negative and
clamps it to zero, and a second sweep reproduces both.

\textbf{A bug worth recording.} An early \texttt{maxViolation} computed
$\max_{i}(a_{i}^{T}x-b_{i})$ but initialised the running maximum to $0$ instead of $-\infty$, so
it could never return a negative value and \emph{strict} feasibility became indistinguishable
from boundary feasibility. Slater's condition was reported as failing for problems that plainly
satisfy it, and the barrier line search --- which accepts a step only when the violation is
negative --- rejected every candidate. The symptom was misleading because the dual path was
unaffected and still produced the correct $\xstar$; only the independent cross-check had silently
stopped running. This is an argument for cross-validation as a design principle rather than a
finishing touch.

\subsection{Plotting engine}

\texttt{PlotCanvas} is a small plotting library written directly on the natID
\texttt{gui::Canvas} primitives, because the SDK installation was missing
\texttt{gui/plot/Plot.h} and \texttt{natPlot} could therefore not be compiled against; writing
the renderer directly also gave the control needed for filled contour fields and envelope
overlays. Data coordinates map affinely to device pixels, with an optional equal-aspect mode that
expands whichever range is short --- necessary in Figure~\ref{fig:primal}, where perpendicularity
of $\nabla f_{0}$ to the contours would otherwise be visually false. Axis ticks use the
$1$--$2$--$5$ decade rule, and label precision is derived from the chosen step.

Scalar fields are sampled on a regular grid and drawn as filled contour bands with curves
overlaid by \textbf{marching squares}, locating edge crossings by linear interpolation. Band
spacing is deliberately non-uniform --- quadratic for the objective field to concentrate bands
near the minimum, and its complement for the dual field to concentrate them near the summit;
uniform spacing wastes almost all bands on the steep outer region. Feasible sets are computed by
\textbf{Sutherland--Hodgman clipping} of a bounding rectangle against each half-plane, and shaded
by $45^{\circ}$ \textbf{hatching} (scanlines in a rotated frame, clipped by Liang--Barsky).
Hatching rather than a solid fill is not decoration: a filled polygon would hide the contours
underneath, and their relation to the feasible set is the whole content of the figure.

\subsection{Interface, export, and safeguards}

Six tabs sit on a \texttt{gui::StandardTabView}; parameter panels use \texttt{gui::NumericEdit}
controls bound to model fields, with a guard flag suppressing the change callback while controls
are populated from the model. Every figure here was produced in batch:
\texttt{-export=<folder>} starts a timer that brings one tab on screen per tick and exports the
previous tab's canvases through \texttt{gui::Canvas::exportToPDF}. Stepping on a timer rather
than looping is what makes it work --- a canvas must be laid out before it has a size to render
into. The canvas palette is chosen from the \emph{rendering backend} rather than the system
theme: on screen it follows dark or light mode, but any other backend always uses the light
palette, because a page is white. Without this, a dark-mode foreground produced white-on-white
PDFs in which the plot lines were visible but every label was invisible.

Because the interface allows arbitrary edits, each solver degrades explicitly: an indefinite $P$
stops the solve and the interface explains that the inner minimisation is unbounded below; a
failed phase I skips the barrier method and reports Slater as unsatisfied while the dual still
provides a valid bound; degenerate rows are skipped by the coordinate ascent; and non-finite
perturbed values are skipped by the hull construction.

% ================================================================= 7
\section{Numerical verification}\label{sec:verify}

All figures below come from the standalone console harness compiled against \texttt{core/} alone.

\begin{center}
\begin{tabular}{llr}
\toprule
\textbf{Check} & \textbf{Measured as} & \textbf{Residual}\\
\midrule
KKT stationarity \eqref{eq:stat} & $\|P\xstar+q+A^{T}\lamstar\|_{\infty}$
  & $5.55\times 10^{-16}$\\
Complementary slackness \eqref{eq:cs} & $\max_{i}|\lambda_{i}^{\star}(b_{i}-a_{i}^{T}\xstar)|$
  & $0.0$\\
Primal feasibility & $\max_{i}(a_{i}^{T}\xstar-b_{i})_{+}$ & $0.0$\\
Dual feasibility & $\max_{i}(-\lambda_{i}^{\star})_{+}$ & $0.0$\\
Two independent primal solves, A & dual route vs.\ log barrier & $1.6\times 10^{-11}$\\
Shadow price, A constraint 1 & central difference vs.\ $-\lambda_{1}^{\star}$
  & $5.6\times 10^{-12}$\\
Shadow price, A constraint 2 & central difference vs.\ $-\lambda_{2}^{\star}$ & $0.0$\\
Shadow price, B budget & central difference vs.\ $-\lamstar$ & $2.4\times 10^{-11}$\\
Convex envelope, C & hull at $u=0$ vs.\ $\dstar$ & $7.6\times 10^{-4}$\\
\bottomrule
\end{tabular}
\end{center}

\vspace{0.4cm}

\begin{center}
\begin{tabular}{lrrrc}
\toprule
\textbf{Instance} & $\pstar$ & $\dstar$ & \textbf{gap} & \textbf{strong duality}\\
\midrule
A --- convex QP   & $-4.250000$ & $-4.250000$ & $0.0$      & holds\\
B --- allocation  & $-8.377168$ & $-8.377168$ & $0.0$      & holds\\
C --- non-convex  & $-7.289309$ & $-9.000000$ & $1.710691$ & \textbf{fails}\\
\bottomrule
\end{tabular}
\end{center}

Weak duality $\dstar\le\pstar$ holds in all three rows, as Theorem~\ref{thm:weak} requires
without exception. Strong duality holds exactly where Slater's hypotheses are met, and fails
where convexity was removed.

% ================================================================= 8
\section{Conclusions}

The dual function was derived in closed form for the QP \eqref{eq:gqp} and the allocation problem
\eqref{eq:waterfill}, and by global minimisation for the non-convex instance \eqref{eq:gnc}; in
all three it came out concave, as Proposition~\ref{prop:concave} requires. Weak duality held
everywhere, consistent with a proof that never invokes convexity. Strong duality held exactly
where Slater's hypotheses were satisfied, and the non-convex instance --- which has strictly
feasible points but a non-convex objective --- produced a gap of $1.7107$, isolating convexity as
the hypothesis doing the work. The identity
$\partial\pstar/\partial u_{i}=-\lambda_{i}^{\star}$ was confirmed to $10^{-11}$ or better in all
three tested cases, including the degenerate one where an inactive constraint correctly priced at
zero, and the identity $\dstar=(\conv p)(0)$ was verified directly against the computed lower
convex hull.

Two methodological points stand apart from the mathematics. Solving each convex primal
\textbf{twice by unrelated methods} turned the duality gap from a tautology into a test --- and
that test caught a real defect that the dual path alone would have concealed. Keeping the
numerical core free of any interface dependency made it testable in isolation, which is where
every result in Section~\ref{sec:verify} was produced.

\textbf{Possible extensions.} Equality constraints with free-sign multipliers; dualising the
non-negativity constraints in Instance B to obtain the full $(n+1)$-dimensional dual; an
active-set or interior-point QP solver for larger $m$; and a search on the constraint level to
locate exactly where the active set changes in the sensitivity curve.

% ================================================================= 9
\begin{thebibliography}{9}

\bibitem{boyd} S.~Boyd and L.~Vandenberghe, \emph{Convex Optimization}, Cambridge University
Press, 2004 --- Ch.~5 (duality), Ch.~11 (interior-point methods).

\bibitem{nocedal} J.~Nocedal and S.~J.~Wright, \emph{Numerical Optimization}, 2nd ed., Springer,
2006 --- Ch.~12, 16.

\bibitem{rockafellar} R.~T.~Rockafellar, \emph{Convex Analysis}, Princeton University Press,
1970 --- conjugacy and perturbation functions.

\bibitem{bertsekas} D.~P.~Bertsekas, \emph{Nonlinear Programming}, 3rd ed., Athena Scientific,
2016 --- duality gaps in non-convex problems.

\bibitem{natid} natID framework documentation and SDK examples, Faculty of Electrical
Engineering, University of Sarajevo.

\end{thebibliography}

\end{document}
